\documentclass[aps,prl,showpacs,amsmath,onecolumn,superscriptaddress,floatfix,nofootinbib,notitlepage]{revtex4-2}

\usepackage{bm,color,xcolor,amsfonts,amsmath,amssymb,epsfig,esint,orcidlink,slashed,braket}

\usepackage[utf8]{inputenc}

\usepackage[T2A]{fontenc}
\usepackage{bbm}
\usepackage{xspace}
\usepackage{enumitem}
\usepackage{tikz}
\usepackage{lipsum}
\usetikzlibrary{calc}
\usetikzlibrary{positioning}
\usepackage[german]{varioref}
\RequirePackage[compat=1.1.0]{tikz-feynman}
\RequirePackage[compat=1.1.0]{tikz-feynhand}
\usepackage{float}
\usepackage{commath}
\usepackage{booktabs}
\usepackage[normalem]{ulem}

\allowdisplaybreaks

\newcommand{\subsubsubsection}[1]{\paragraph{#1} ~ \vspace{1.7ex plus 0.2ex}}

\newenvironment{todo}%
{\large \textsc{Todo:} \normalsize \begin{itemize} \color{blue}}%
{\end{itemize}}

\newcommand{\red}[1]{{\color{red}#1}}
\newcommand{\blue}[1]{{\color{blue}#1}}
\DeclareRobustCommand{\magenta}[1]{{\color{magenta}#1}}
\newcommand{\fkcomment}[1]{{\color{blue}[FK: #1]}}
\newcommand{\chhcomment}[1]{{\color{red}[CH: #1]}}
\newcommand{\ancomment}[1]{{\color{blue}[AN: #1]}}
\newcommand{\fk}[1]{{\color{blue}#1}}
\newcommand{\ch}[1]{{\color{green}#1}}
\newcommand{\an}[1]{{\color{blue}#1}}
\newcommand{\pp}[1]{{\color{orange}[PP: #1]}}

\newcommand{\bbnP}{[\bar{b}b]_1(nP)}
\newcommand{\bbarb}{[\bar{b}b]_1}
%\newcommand{\tauEB}{\hat{\tau}^{-1}(-E_B)}
\newcommand{\tauEB}{\hat{\cal J}}
\newcommand{\tauinv}{\hat\tau^{-1}}
\newcommand{\dtauinvEB}{\Bigl[\partial_E\hat\tau^{-1}\Bigr]_{E=-E_B}}
\newcommand{\psibb}{\Psi^{(\bar{b}b)}}
\newcommand{\Mbb}{M^{(\bar{b}b)}}

\setcounter{secnumdepth}{3}

\newcommand{\GeV}{\mathrm{GeV}}
\newcommand{\MeV}{\mathrm{MeV}}
\newcommand{\vep}{{\bm p}}
\newcommand{\vek}{{\bm k}}
\newcommand{\veq}{{\bm q}}
\newcommand{\ver}{{\bm r}}
\newcommand{\vex}{{\bm x}}
\newcommand{\veA}{{\bm A}}
\newcommand{\vetau}{{\bm \tau}}
\definecolor{green}{RGB}{44,160,44}
\definecolor{omeRevision}{RGB}{0,105,105}
\newcommand{\be}{\begin{equation}}
\newcommand{\ee}{\end{equation}}
\newcommand{\bea}{\begin{eqnarray}}
\newcommand{\eea}{\end{eqnarray}}
\newcommand{\beas}{\begin{eqnarray*}}
\newcommand{\eeas}{\end{eqnarray*}}
\newcommand{\ds}{\displaystyle}
\newcommand{\B}{\color{blue}}
\newcommand{\R}{\color{red}}
\newcommand{\G}{\color{green}}
\newcommand{\veL}{{\bm L}}
\newcommand{\veS}{{\bm S}}
\newcommand{\ven}{{\bm n}}

\newcommand{\nn}{\notag}
\newcommand{\dk}{\frac{d^3k}{(2\pi)^3}}
\newcommand{\ddq}{\frac{d^3q}{(2\pi)^3}}
\newcommand{\ddqpr}{\frac{d^3q'}{(2\pi)^3}}
\newcommand{\dqr}{\frac{q^2dq}{2\pi^2}}
\newcommand{\dqrpr}{\frac{q^{\prime2}dq'}{2\pi^2}}
\newcommand{\ddp}{\frac{d^3p}{(2\pi)^3}}
\newcommand{\ddpp}{\frac{d^3p'}{(2\pi)^3}}
\newcommand{\intp}{\int\frac{d^3p}{(2\pi)^3}}
\newcommand{\intk}{\int\frac{q^2dq}{(2\pi)^3}}
\newcommand{\bb}{\bar{b}b}
\newcommand{\qq}{q\bar{q}}
\newcommand{\QQ}{\bar{b}b}
\newcommand{\gpi}{g_b}
\renewcommand{\vec}{\bm}
\newcommand{\no}{\notag\\}
% \newcommand{\Hoo}{\hat E_1}
\newcommand{\Hoo}{H_{\rm MM}}
\newcommand{\Hot}{V_{ls}}
\newcommand{\Htt}{V_{s}}
\newcommand{\mG}{m_\phi}
\newcommand{\EG}{E_\phi}
\newcommand{\CSOME}{C_S^{\rm (OME)}}
\newcommand{\COOME}{C_O^{\rm (OME)}}
\newcommand{\CSQM}{C_S^{(\QQ)}}
\newcommand{\COQM}{C_O^{(\QQ)}}
\newcommand{\CSOQM}{C_{SO}^{(\QQ)}}
\newcommand{\En}{E_n}
\newcommand{\Enzero}{E_n^{(0)}}
\newcommand{\Eizero}{E_i^{(0)}}
\newcommand{\VR}{V_R}
\newcommand{\gc}{g_c}
\newcommand{\X}{X(3872)}
\newcommand{\Xb}{X_b}
\newcommand{\gev}{GeV}
\newcommand{\mev}{MeV}
\newcommand{\Yb}{\Upsilon(10753)}
\newcommand{\Yfour}{\Upsilon(10580)}
\newcommand{\Yfive}{\Upsilon(10860)}
\newcommand{\Ysix}{\Upsilon(11020)}
\newcommand{\Dszero}{D^*_{s0}(2317)}
\newcommand{\veep}{{\bm \epsilon}}
\newcommand{\veepst}{{\bm \epsilon}^{\prime*}}
\newcommand{\tbd}[1]{{\color{red}[TODO: #1]}}
\newcommand{\unit}[1]{~(#1)}
\newcommand{\BOchiEFT}{BO$\chi$EFT\xspace}

\newcommand{\VQQ}{V^{(\bar{b}b)}}
\newcommand{\VOME}{V^{\rm (OME)}}
\newcommand{\calH}{{\cal H}}
\newcommand{\hatH}{H}
\newcommand{\tilH}{\tilde{{\cal H}}}


\newcommand{\rub}{\affiliation{Institut f\"ur Theoretische Physik II, Ruhr-Universit\"at Bochum, D-44780 Bochum, Germany }}

\newcommand{\bonn}{\affiliation{Helmholtz-Institut f\"ur Strahlen- und Kernphysik and Bethe Center for Theoretical Physics,\\ Universit\"at Bonn, D-53115 Bonn, Germany}}

\newcommand{\fzj}{\affiliation{Institute for Advanced Simulation
and Cluster of Excellence ``Color meets Flavor'', Forschungszentrum J\"ulich, D-52425 J\"ulich, Germany}}

\newcommand{\itp}{\affiliation{Institute of Theoretical Physics, Chinese Academy of Sciences, Beijing 100190, China}}

\newcommand{\ucas}{\affiliation{School of Physical Sciences, University of Chinese Academy of Sciences, Beijing 100049, China}}

\newcommand{\lis}{\affiliation{CeFEMA, Center of Physics and Engineering of Advanced Materials, Instituto Superior T{\'e}cnico, Avenida Rovisco Pais 1, 1049-001 Lisboa, Portugal}}

\newcommand{\tum}{\affiliation{
Technical University of Munich, TUM School of Natural Sciences, Physics Department, 85748 Garching, Germany}}

\newcommand{\scnt}{\affiliation{Southern Center for Nuclear-Science Theory (SCNT), Institute of Modern Physics,\\ Chinese Academy of Sciences, Huizhou 516000, China}}

\newcommand{\hiskp}{\affiliation{Helmholtz-Institut f\"ur Strahlen- und Kernphysik% %and Bethe Center for Theoretical Physics
, Universit\"at Bonn, 53115 Bonn, Germany}}

\newcommand{\hohai}{\affiliation{College of Mechanics and Engineering Science, Hohai University, Nanjing 211100, China}}


\graphicspath{{Figs.dir/}}

\begin{document}

\title{Follow-ups}

%\author{Yu Zhang\orcidlink{0009-0007-2951-6536}}\email{zhangyu@itp.ac.cn}
%\itp\ucas

% \author{Feng-Kun~Guo\orcidlink{0000-0002-2919-2064}}
% \email{fkguo@itp.ac.cn}
% \itp\ucas\scnt
%
% \author{Christoph~Hanhart\orcidlink{0000-0002-3509-2473}}\email{c.hanhart@fz-juelich.de}
% \fzj
%
% \author{Alexey Nefediev\orcidlink{0000-0002-9988-9430}}\email{nefediev@hiskp.uni-bonn.de}
% \hiskp

\maketitle

\tableofcontents

\section{$\Xb$ spin-2 partner}

We have the interaction in the $b$-sector fully fixed in Ref.~\cite{Zhang:2026iyr}. With these parameters, we proceed to the $J^{PC}=2^{++}$ sector and make prediction for the spin-2 partner of the $\Xb$. This is possible since the contact term $C$ is the same in the $1^{++}$ and $2^{++}$ sectors and
\be
g[(\bar{b}b)_2\mbox{-}B^*\bar{B}^*(^5S_2)]=
g[(\bar{b}b)_1\mbox{-}B\bar{B}^*(^3S_1)]=\frac{2}{\sqrt{3}}\sqrt{2}g_l
\ee
since the recoupling coefficient in the $2^{++}$ channel is  $-1/\sqrt{3}$ (see Table~2 of Ref.~\cite{Bruschini:2023zkb}).
A complete information on the coupled open-bottom channels, partial waves, and projectors can be found in Ref.~\cite{Baru:2019xnh}.

In the $2^{++}$ channel with the threshold at 10.650~GeV, we have three generic bottomonium states with the masses
\be
\begin{split}
M[\chi_{b2}(1P)]&=9912.21~\mbox{MeV},\\
M[\chi_{b2}(2P)]&=10268.65~\mbox{MeV},\\
M[\chi_{b2}(3P)]&=10524.0~\mbox{MeV}.
\end{split}
\ee

It should be checked how well their masses are reproduced with the parameters fixed in the $1^{++}$ sector. After that, a molecular pole representing the $X_{b2}$ can be searched for in the two scenarious discussed in our paper \cite{Zhang:2026iyr}:
(i) the scenario with $\Delta M=0$ and (ii) the scenario with the $\Xb$ binding energy fixed to 1~MeV below the $B\bar{B}^*$ threshold.

The six-channel OME potential and all 21 independent partial-wave entries are derived in notes/04-ome-2pp.md and presented in the appendix of paper/main.tex.

\subsection{OME in $2^{++}$}
\begin{center}
    
    \begin{tabular*}{0.5\textwidth}{@{\extracolsep{\fill}}lccc@{}}    \toprule
    pair & $S$ & $L$ & $^{2S+1}L_J$\\
    \midrule
    $B\bar{B}$ & 0   &  2 & $^1D_2$\\
    $B\bar{B}^*$, $C$-even & 1 & 2 & $^3D_2$\\
    $B^*\bar{B}^*$ & 0 & 2 & $^1D_2$\\
    $B^* \bar{B}^*$ & 2 & 0,2,4 & $^5S_2$, $^5D_2$, $^5G_2$\\
    \bottomrule
    \end{tabular*}
\end{center}

All the possible $2^{++}$ channels
\be
B^*\bar{B}^*(^5S_2),\quad B^*\bar{B}^*(^1D_2),\quad B^*\bar{B}^*(^5D_2),\quad B^*\bar{B}^*(^5G_2),\quad B\bar{B}^*(^3D_2),\quad B\bar{B}(^1D_2).
\ee

Since $BB\phi$ vertex violates parity and is absent in lagrangian, so the interactions:
\be
V_{B\bar{B}->B\bar{B}}, \quad V_{B\bar{B} -> B\bar{B}^*},\quad V_{B\bar{B}^* -> B\bar{B}}
\ee
all vanish and the only new transition is $B\bar{B} \leftrightarrow B^*\bar{B}^*$.
\be
\begin{split}
V_{B\bar{B}\rightarrow B^*\bar{B}^*}^{(\text{OME})}(\vep',\vep; m_\phi) = \frac{g_b^2}{4F_\pi^2}(\veepst_1 \cdot \veq)(\veepst_2 \cdot \veq)\left[2D_3(\vep',\vep;m_\phi)\right]\\
V_{B^*\bar{B}^*\rightarrow B\bar{B}}^{(\text{OME})}(\vep',\vep; m_\phi) = \frac{g_b^2}{4F_\pi^2}(\veep_1 \cdot \veq)(\veep_2 \cdot \veq)\left[2D_4(\vep',\vep;m_\phi)\right]
\end{split}
\ee

The partial-wave projections are given in the appendix of paper/main.tex.

\section{$\X$ and its spin-2 partner}

We have three known $J^{PC}=1^{++}$ states in the spectrum of charmonium:
\begin{itemize}
\item $\chi_{c1}(1P)$ as a generic charmonium with
\be
M[\chi_{c1}(1P)]=3510.67~\mbox{MeV},\quad \Gamma[\chi_{c1}(1P)]=0.82\pm 0.04~\mbox{MeV}.
\ee
\item $\X$ as a pure molecular state with
\be
E_B[\X]\simeq 160~\mbox{keV}.
\ee
\item $\chi_{c1}(4010)$ as candidate to generic $\chi_{c1}(2P)$ state but lying above the open-charm threshold $D\bar{D}^*$, with \cite{LHCb:2024vfz}
\be
M[\chi_{c1}(4010)]=4012.5\,_{-3.9}^{+3.6}{}\,_{-3.7}^{+4.1}~\mbox{MeV},\quad \Gamma[\chi_{c1}(4010)]=62.7\,_{-6.4}^{+7.0}{}\,_{-6.6}^{+6.4}~\mbox{MeV}.
\ee
\end{itemize}

Thus, we first repeat the calculation of the paper \cite{Zhang:2026iyr} in the charm sector:
\begin{itemize}
\item The $c$-quark mass $m_c$ is fixed to reproduce the splitting between the $\chi_{c1}(1P)$ and $\chi_{c1}(4010)$ (taken as generic $\chi_{c1}(2P)$ meson dressed through the interaction with the open-charm channel),
\be
m_c=1.328~\mbox{GeV},
\ee
which is consistent with natural values of the $c$-quark mass adopted in various quark model calculations.
\item The parameter $V_0$ is fixed to reproduce the exact position of the $\chi_{c1}$ states in the spectrum,
\be
V_0=-36~\mbox{MeV}.
\ee
\item The contact term $C$ is fixed to have a shallow bound molecular state $\X$.
\end{itemize}
The LSE for the coupled $D\bar{D}$, $D\bar{D}^*$, and $D^*\bar{D}^*$ channels with all partial waves are identical to those in the bottom sector. Thus we simply repeat the calculation for the $X_{b2}$ and make prediction for the $\X$ spin-2 partner $X_2$.

Preliminary work needed:
\begin{itemize}
\item The pole search procedure needs to be generalised to above-threshold poles since $\chi_{c1}(4010)$ lies above the $D\bar{D}^*$ threshold. The poles lies in the complex plane away from the real axis and has a sizable imeginary part around 60~MeV.
\end{itemize}

\section{CMS line shape analysis}

\begin{thebibliography}{99}
%\cite{Zhang:2026iyr}
\bibitem{Zhang:2026iyr}
Y.~Zhang, F.~K.~Guo, C.~Hanhart, A.~Nefediev and P.~P{\"u}tz,
%``Born-Oppenheimer chiral effective field theory and its application to $X_b$,''
[arXiv:2610.19066 [hep-ph]].
%0 citations counted in INSPIRE as of 19 Oct 2026
%\cite{Baru:2019xnh}
%\cite{Bruschini:2023zkb}
\bibitem{Bruschini:2023zkb}
R.~Bruschini,
%``Heavy-quark spin symmetry breaking in the Born-Oppenheimer approximation,''
JHEP \textbf{08}, 219 (2023)
doi:10.1007/JHEP08(2023)219
[arXiv:2303.17533 [hep-ph]].
%23 citations counted in INSPIRE as of 19 Oct 2026
\bibitem{Baru:2019xnh}
V.~Baru, E.~Epelbaum, A.~A.~Filin, C.~Hanhart, A.~V.~Nefediev and Q.~Wang,
%``Spin partners $W_{bJ}$ from the line shapes of the $Z_b(10610)$ and $Z_b(10650)$,''
Phys. Rev. D \textbf{99}, no.9, 094013 (2019)
doi:10.1103/PhysRevD.99.094013
[arXiv:1901.10319 [hep-ph]].
%55 citations counted in INSPIRE as of 19 Oct 2026

%\cite{LHCb:2024vfz}
\bibitem{LHCb:2024vfz}
R.~Aaij \textit{et al.} [LHCb],
%``Observation of New Charmonium or Charmoniumlike States in B+{\textrightarrow}D*{\ensuremath{\pm}}D{\ensuremath{\mp}}K+ Decays,''
Phys. Rev. Lett. \textbf{133}, no.13, 131902 (2024)
doi:10.1103/PhysRevLett.133.131902
[arXiv:2406.03156 [hep-ex]].
%72 citations counted in INSPIRE as of 21 Oct 2026

\end{thebibliography}


\end{document}
